\documentclass[11pt]{article}
\usepackage[T1]{fontenc}
\usepackage[utf8]{inputenc}
\usepackage{lmodern}
\usepackage[a4paper,margin=28mm]{geometry}
\usepackage{microtype}
\usepackage{amsmath,amssymb,amsthm,mathtools}
\usepackage{enumitem}
\usepackage{xcolor}
\definecolor{linkblue}{rgb}{0.12,0.22,0.42}
\usepackage[colorlinks=true,allcolors=linkblue]{hyperref}
\usepackage[capitalise,nameinlink,noabbrev]{cleveref}
\hypersetup{pdftitle={Electromagnetic tails and the Kerr--Newman Penrose inequality},pdfauthor={OpenAI}}
\numberwithin{equation}{section}
\newtheorem{theorem}{Theorem}[section]
\newtheorem{lemma}[theorem]{Lemma}
\newtheorem{proposition}[theorem]{Proposition}
\newtheorem{corollary}[theorem]{Corollary}
\theoremstyle{definition}
\newtheorem{definition}[theorem]{Definition}
\newtheorem{remark}[theorem]{Remark}
\newcommand{\R}{\mathbb R}
\newcommand{\Sph}{\mathbb S}
\newcommand{\dd}{\mathop{}\!\mathrm d}
\DeclareMathOperator{\tr}{tr}
\DeclareMathOperator{\id}{id}
\DeclareMathOperator{\supp}{supp}
\DeclareMathOperator{\Area}{Area}
\DeclareMathOperator{\Vol}{Vol}
\DeclareMathOperator{\End}{End}
\newcommand{\norm}[1]{\left\lVert#1\right\rVert}
\newcommand{\abs}[1]{\left\lvert#1\right\rvert}
\setlist[enumerate]{label=\textup{(\roman*)},leftmargin=*,itemsep=3pt}
\setlist[itemize]{leftmargin=*,itemsep=3pt}
\allowdisplaybreaks[1]

\title{Electromagnetic tails and the\texorpdfstring{\\}{}Kerr--Newman Penrose inequality}
\author{OpenAI}
\date{October 5, 2026}
\begin{document}
\maketitle
\begin{abstract}
We construct smooth axisymmetric electrovacuum exteriors that violate the
Kerr--Newman Penrose inequality when \(J\) is the bare gravitational ADM
angular momentum and the electromagnetic fields have only \(O(r^{-2})\)
decay, allowing angularly varying leading tails. The examples have zero
total electric and magnetic charge even though both electromagnetic fields
are nonzero. They have a connected outermost, outer-area-minimizing future
marginally outer trapped boundary and lie strictly on the physical area
branch. We also obtain equality examples whose original data admit no
Kerr--Newman spacelike realization with the corresponding mass, angular
momentum, and charges. These counterexamples do not refute formulations
using conserved total angular momentum with its electromagnetic correction
or stronger Coulomb asymptotics.
\end{abstract}
\section{Introduction}\label{sec:introduction}

The Kerr--Newman mass formula suggests the Kerr--Newman Penrose inequality
\begin{equation}\label{eq:kn-inequality}
 m^2\ge \frac{A}{16\pi}+\frac{Q^2}{2}
             +\frac{\pi(Q^4+4J^2)}{A},
 \qquad A\ge4\pi\sqrt{Q^4+4J^2},
\end{equation}
where \(A\) is a suitable enclosing area, \(m\) is mass, \(Q\) is total
charge, and \(J\) is angular momentum. Its motivation comes from the
gravitational-collapse argument underlying Penrose's original proposal
\cite{Penrose1973}; the area branch in \eqref{eq:kn-inequality} is the branch
on which the right side increases with \(A\). The formulation using minimum
enclosing area, and the distinction between the corresponding upper and
lower area bounds, are discussed in \cite{DKWY2013}.

In the presence of an electromagnetic field, specifying \(J\) is essential.
The gravitational angular-momentum flux through a surface need not be
conserved. For the globally defined invariant magnetic potential used
below, adding the corresponding electromagnetic boundary term gives a
conserved flux under the electrovacuum constraints.
Khuri--Weinstein include this term in their definition at infinity and
state that its disappearance requires appropriate asymptotic conditions
\cite[Equation~(1.12) and the following paragraph]{KW2016}.
The identification argument of Dain--Khuri--Weinstein--Yamada uses
Coulomb expansions for the electric and magnetic fields
\cite[Equations~(2.20) and~(2.22)]{DKWY2013}.
The weaker bounds \(\mathcal E,\mathcal B=O(r^{-2})\) alone do not impose
those leading terms.

We give a negative resolution of \eqref{eq:kn-inequality} for the precise
formulation below: \(J\) is the gravitational ADM flux at infinity, and the
electromagnetic assumptions are only the displayed decay and constraint
conditions. Our examples have a connected outermost future marginally
outer trapped boundary, minimize area among all enclosing cuts, and lie
strictly on the physical area branch. They also disprove the associated
nondegenerate Kerr--Newman equality implication. Stronger asymptotic
conditions, or a definition of angular momentum retaining its
electromagnetic term, lead to different assertions; no conclusion about
those assertions follows from these examples.

\subsection{Conventions and the class of exteriors}\label{sec:conventions}

We work in three spatial dimensions with \(G=c=1\) and zero cosmological
constant. On an oriented Riemannian three-manifold \((\Omega,g)\), the cross
product is defined by
\[
 (X\times_gY)^\flat(Z)=\dd V_g(X,Y,Z).
\]
Let \(K\) be a symmetric covariant two-tensor, \(\tau=\tr_gK\), and let
\(\mathcal E,\mathcal B\) be electric and magnetic vector fields. We use
the source-free electrovacuum constraint equations
\begin{equation}\label{eq:constraints}
 \begin{split}
 R_g+\tau^2-\abs{K}_g^2
      &=2\bigl(\abs{\mathcal E}_g^2+\abs{\mathcal B}_g^2\bigr),\\
 \operatorname{div}_g(K-\tau g)
      &=2(\mathcal E\times_g\mathcal B)^\flat,\\
 \operatorname{div}_g\mathcal E
      &=\operatorname{div}_g\mathcal B=0.
 \end{split}
\end{equation}
Thus the total constraint densities are
\[
 \mu=\frac{\abs{\mathcal E}_g^2+\abs{\mathcal B}_g^2}{8\pi},
 \qquad J_{\mathrm{con}}
       =\frac{(\mathcal E\times_g\mathcal B)^\flat}{4\pi}.
\]
The covector \(J_{\mathrm{con}}\) is distinct from the scalar angular
momentum \(J\).

All data are smooth, including at the compact boundary. The exterior is
connected, complete as a metric space with its boundary included, and
diffeomorphic to the complement of an open ball in \(\R^3\). There is one
asymptotically flat end and a smooth effective \(U(1)\) action preserving
\(g,K,\mathcal E,\mathcal B\) and the boundary \(S\). Its generator \(\eta\)
has period \(2\pi\) and agrees on the end with the standard oriented
rotation about the third coordinate axis.

In asymptotic Euclidean coordinates \(x\), write \(r=\abs{x}\) and assume
\begin{equation}\label{eq:falloff}
 g_{ij}-\delta_{ij}=O_4(r^{-1}),\qquad
 K_{ij}=O_3(r^{-3}),\qquad
 \mathcal E^i,\mathcal B^i=O_3(r^{-2}).
\end{equation}
Here \(O_k(r^{-a})\) bounds each Cartesian derivative of order
\(j\le k\) by \(Cr^{-a-j}\). We require \(\mu\) and
\(\abs{J_{\mathrm{con}}}_g\) to be integrable. On a coordinate sphere
\(S_R\), let \(n_\delta\) be its Euclidean outward unit normal and \(\nu_R\)
its \(g\)-unit outward normal. The asymptotic quantities are
\begin{align}
 E_{\mathrm{ADM}}
   &=\frac1{16\pi}\lim_{R\to\infty}
       \int_{S_R}(\partial_jg_{ij}-\partial_ig_{jj})
           n_\delta^i\,\dd A_\delta,\label{eq:adm-energy}\\
 (P_{\mathrm{ADM}})_i
   &=\frac1{8\pi}\lim_{R\to\infty}
       \int_{S_R}(K_{ij}-\tau g_{ij})n_\delta^j\,\dd A_\delta,
       \label{eq:adm-momentum}\\
 J&=\frac1{8\pi}\lim_{R\to\infty}
       \int_{S_R}(K_{ij}-\tau g_{ij})\nu_R^i\eta^j\,\dd A_g,
       \label{eq:adm-angular}\\
 Q_e&=\frac1{4\pi}\lim_{R\to\infty}
       \int_{S_R}g(\mathcal E,\nu_R)\,\dd A_g,
 \qquad
 Q_b=\frac1{4\pi}\lim_{R\to\infty}
       \int_{S_R}g(\mathcal B,\nu_R)\,\dd A_g.
       \label{eq:charges-definition}
\end{align}
All these limits are required to exist and be finite. We impose
\(P_{\mathrm{ADM}}=0\) and \(E_{\mathrm{ADM}}>0\), and set
\(m=E_{\mathrm{ADM}}\), \(Q=\sqrt{Q_e^2+Q_b^2}\).
In particular \eqref{eq:adm-angular} has no electromagnetic surface term.

For a two-sided surface with unit normal \(\nu\) pointing toward the end,
put
\begin{equation}\label{eq:expansion}
 H_\Sigma=\operatorname{div}_\Sigma\nu,
 \qquad \tr_\Sigma K=(g^{ij}-\nu^i\nu^j)K_{ij},
 \qquad \theta_+(\Sigma)=H_\Sigma+\tr_\Sigma K.
\end{equation}
Thus outward Euclidean spheres have positive mean curvature. The boundary
is a future marginally outer trapped surface, or future MOTS, if
\(\theta_+(S)=0\). We require this condition and the following precise
outermostness condition: there is no compact smooth embedded enclosing
surface entirely in the interior, possibly disconnected, for which
\(\theta_+\le0\) everywhere. No condition on the inward expansion is imposed.

An \emph{enclosing cut} is the compact intrinsic boundary
\(\Gamma=\partial D\) of a connected smooth codimension-zero submanifold
\(D\subset\Omega\), closed as a subset of \(\Omega\), whose interior lies
in \(\operatorname{int}\Omega\) and which contains the entire sufficiently
distant part of the end. The full intrinsic boundary is counted, including
any portion coinciding with \(S\). The cut must be smooth, embedded and
two-sided; it may be disconnected and need not be symmetric or trapped.
In particular \(D=\Omega\) gives the cut \(S\). The required area condition is
\begin{equation}\label{eq:outer-area-condition}
 \Area_g(\Gamma)\ge A:=\Area_g(S)>0
 \qquad\text{for every enclosing cut }\Gamma.
\end{equation}
Thus \(A\) is exactly the minimum enclosing area.

\subsection{Main result}\label{sec:main-result}

\begin{theorem}\label{thm:main}
In the class of smooth axisymmetric source-free electrovacuum exteriors
specified in \Cref{sec:conventions}, both of the following occur.
\begin{enumerate}
 \item There are data with
 \[
 Q_e=Q_b=0,\qquad A>8\pi\abs J,
 \qquad m^2<\frac{A}{16\pi}+\frac{4\pi J^2}{A}.
 \]
 \item There are data with
 \[
 Q_e=Q_b=0,\qquad A>8\pi\abs J,
 \qquad m^2=\frac{A}{16\pi}+\frac{4\pi J^2}{A},
 \]
 for which no smooth spacelike embedding into the Kerr--Newman spacetime
 with parameters \((m,J,Q_e,Q_b)\) induces the given metric, second
 fundamental form, and electromagnetic fields.
\end{enumerate}
All these examples are maximal, \(\tr_gK=0\), on
\(\Omega=\{x\in\R^3:\abs{x}\ge1\}\), with connected boundary
\(S=\{r=1\}\). Both electromagnetic fields are nonzero.
\end{theorem}

The first conclusion disproves \eqref{eq:kn-inequality} under
\eqref{eq:falloff} and \eqref{eq:adm-angular}, even with every horizon and
area condition above. The second disproves the nondegenerate equality
implication with a Kerr--Newman realization, including realizations by
nonconstant time graphs. The usual spacetime convention is
\(K(X,Y)=\overline g(\overline\nabla_X n,Y)\) for the future unit normal
\(n\); the induced fields are
\(\mathcal E^\flat=(\iota_nF)|_{T\Omega}\) and
\(\mathcal B^\flat=\star_g(F|_{T\Omega})\).
The obstruction in the second conclusion holds for every spacelike slice,
so it also excludes a realization extending smoothly to a cross-section
of the future horizon or to the bifurcation sphere.

The proof is independent of any neutral or charged Penrose inequality or
rigidity theorem. It gives a two-parameter family, with amplitude
\(s\in[0,1]\) and a transition radius \(L\), such that
\begin{equation}\label{eq:family-summary}
 J=-\frac{4s^2}{15},\qquad
 m=2+O(s^2/L),\qquad
 A=64\pi+O(s^2/L^2).
\end{equation}
The electromagnetic correction to the angular-momentum flux is
\(+4s^2/15\) at infinity. Angularly varying radial tails supply this term
while their energy tends to zero as \(L\to\infty\). At \(s=1\), the defect
in \eqref{eq:kn-inequality} tends to \(-1/225\). For fixed large \(L\), a
small positive amplitude gives the opposite sign, yielding the equality
example by continuity.

Two features make the construction explicit. A mixed radial--azimuthal
tensor solves the electromagnetic momentum constraint without a vector
elliptic equation. Kelvin reflection then gives a positive Newton-potential
operator with the Robin condition that makes the unit sphere minimal.
The contraction estimate uses the distance of the source from that sphere,
and accommodates a noncompact \(r^{-4}\) energy tail. Finally, a closed
two-form obtained by radial projection proves area-minimization against
all enclosing cuts. These steps are developed in
\Cref{sec:construction,sec:elliptic,sec:geometry};
\Cref{sec:counterexamples} computes the fluxes and proves both conclusions.

\section{Explicit electrovacuum seeds}
\label{sec:construction}

We first solve the momentum and Maxwell divergence constraints explicitly.
The electromagnetic fields begin at a large radius; one scalar equation will
then determine the metric.  Throughout the construction,
\[
 \Omega=\{x\in\R^3:r=|x|\ge1\},\qquad S=\{r=1\},
\]
with the Euclidean metric $\delta$ and standard orientation.  Write
$(r,\theta,\phi)$ for spherical coordinates, with $\theta$ the colatitude and
$\eta=\partial_\phi$ the generator of the standard rotation action.
Fix a smooth function $\chi:\R\to[0,1]$ that is zero on $(-\infty,1]$ and
one on $[2,\infty)$.  For example, one may take
\[
 \chi(t)=\frac{\zeta(t-1)}{\zeta(t-1)+\zeta(2-t)},\qquad
 \zeta(t)=\begin{cases}e^{-1/t},&t>0,\\0,&t\le0.\end{cases}
\]
For $s\in[0,1]$ and $L>4$, put $c(r)=\chi(r/L)$.  Constants $C,C_j$ below
depend only on this fixed cutoff and the indicated derivative order, and
are independent of $s$ and $L$ above a fixed lower threshold.

Define the scalar potentials, vector fields, and symmetric covariant tensor
by
\begin{equation}
\label{eq:seeds}
\begin{aligned}
 f&=sc(r)\sin^2\theta,& h&=sc(r)\sin^2\theta\cos\theta,\\
 \iota_{\widetilde E}\dd V_\delta&=\dd(f\,\dd\phi),&
 \iota_{\widetilde B}\dd V_\delta&=\dd(h\,\dd\phi),\\
 \sigma&=-\frac{s^2c(r)^2\sin^4\theta}{r^2}
       (\dd r\otimes\dd\phi+\dd\phi\otimes\dd r).
\end{aligned}
\end{equation}
Here contraction with the volume form uniquely determines each vector
field.  The coordinate formulas in \eqref{eq:seeds} extend smoothly across
the rotation axis, as \Cref{lem:seeds} verifies.

\begin{lemma}[Seed identities and bounds]
\label{lem:seeds}
The fields $\widetilde E,\widetilde B$ and tensor $\sigma$ are smooth on
$\Omega$, invariant under the standard effective $U(1)$ action, and vanish
on $\{1\le r\le L\}$.  They satisfy
\begin{equation}
\label{eq:seed-constraints}
 \operatorname{div}_\delta\widetilde E
 =\operatorname{div}_\delta\widetilde B=0,\qquad
 \tr_\delta\sigma=0,\qquad
 \operatorname{div}_\delta\sigma
 =2\dd V_\delta(\widetilde E,\widetilde B,\cdot).
\end{equation}
For every Cartesian multi-index $\alpha$ of order $j\ge0$,
\begin{equation}
\label{eq:seed-bounds}
\begin{aligned}
 |\partial^\alpha\widetilde E|_\delta
 +|\partial^\alpha\widetilde B|_\delta
 &\le C_j s r^{-2-j}\mathbf1_{\{r\ge L\}},\\
 |\partial^\alpha\sigma|_\delta
 &\le C_j s^2 r^{-3-j}\mathbf1_{\{r\ge L\}}.
\end{aligned}
\end{equation}
Moreover, $\operatorname{div}_\delta\sigma$ is supported in
$\{L\le r\le2L\}$, with
\[
 |\partial^\alpha(\operatorname{div}_\delta\sigma)|_\delta
 \le C_j s^2r^{-4-j}\mathbf1_{\{L\le r\le2L\}}.
\]
\end{lemma}

\begin{proof}
Let $X=(x^1,x^2,x^3)$, $T=(-x^2,x^1,0)$, and
$w=(x^1)^2+(x^2)^2$, with $\flat$ denoting Euclidean metric dual in the
following formulas.  The primitives and tensor have the Cartesian expressions
\begin{equation}
\label{eq:cartesian-seeds}
\begin{aligned}
 f\,\dd\phi&=\frac{sc}{r^2}(x^1\dd x^2-x^2\dd x^1),\\
 h\,\dd\phi&=\frac{scx^3}{r^3}(x^1\dd x^2-x^2\dd x^1),\\
 \sigma&=-\frac{s^2c^2w}{r^7}
   (X^\flat\otimes T^\flat+T^\flat\otimes X^\flat).
\end{aligned}
\end{equation}
These are smooth on $r\ge1$, with no denominator involving distance to
the axis.  They also show rotation invariance and, since $X\cdot T=0$,
trace-freeness.  The exact flux forms are closed, so
$\dd(\iota_{\widetilde E}\dd V_\delta)=
(\operatorname{div}_\delta\widetilde E)\dd V_\delta=0$, and likewise for
$\widetilde B$.

For the momentum identity we compute off the axis and then use smoothness.
Since $\dd V_\delta=r^2\sin\theta\,\dd r\wedge\dd\theta\wedge\dd\phi$,
the spherical \emph{coordinate} components are
\begin{equation}
\label{eq:field-components}
\begin{aligned}
 \widetilde E^r&=\frac{2sc\cos\theta}{r^2},&
 \widetilde E^\theta&=-\frac{sc'\sin\theta}{r^2},&
 \widetilde E^\phi&=0,\\
 \widetilde B^r&=\frac{sc(3\cos^2\theta-1)}{r^2},&
 \widetilde B^\theta&=-\frac{sc'\sin\theta\cos\theta}{r^2},&
 \widetilde B^\phi&=0.
\end{aligned}
\end{equation}
In particular an orthonormal polar component is $r$ times the displayed
polar component.  The covector $\dd V_\delta(\widetilde E,\widetilde B,\cdot)$
has only an azimuthal component, namely
\begin{equation}
\label{eq:seed-cross-product}
 \dd V_\delta(\widetilde E,\widetilde B,\partial_\phi)
 =\frac{f_rh_\theta-f_\theta h_r}{r^2\sin\theta}
 =-\frac{s^2cc'\sin^4\theta}{r^2}.
\end{equation}
To check all three tensor-divergence components, put
$a=\sigma_{r\phi}=-s^2c^2\sin^4\theta/r^2$ and use the coordinate identity
\[
 (\operatorname{div}_\delta\sigma)_i
 =\frac1{\sqrt{\det\delta}}\partial_j
    \bigl(\sqrt{\det\delta}\,\sigma_i{}^j\bigr)
   -\frac12\sigma^{jk}\partial_i\delta_{jk}.
\]
The last term vanishes: the metric is diagonal in these coordinates,
whereas $\sigma^{jk}$ has only $r\phi$ and $\phi r$ entries.  Thus
\begin{equation}
\label{eq:all-divergence-components}
\begin{aligned}
 (\operatorname{div}_\delta\sigma)_r
  &=\frac1{r^2\sin\theta}\partial_\phi
       \left(\frac a{\sin\theta}\right)=0,\\
 (\operatorname{div}_\delta\sigma)_\theta&=0,\\
 (\operatorname{div}_\delta\sigma)_\phi
  &=\frac1{r^2\sin\theta}\partial_r(r^2\sin\theta\,a)
    =a_r+\frac{2a}{r}
    =-\frac{2s^2cc'\sin^4\theta}{r^2}.
\end{aligned}
\end{equation}
This proves \eqref{eq:seed-constraints}.  The last expression also proves
the stated support property; for $r\ge2L$ both fields are radial and
parallel.

Finally, for $j\ge1$ the Cartesian derivatives of $c$ obey
\[
 |\partial^\alpha c|\le C_j r^{-j}
   \mathbf1_{\{L\le r\le2L\}}.
\]
Indeed, all such derivatives are supported where $r$ is comparable to
$L$, and differentiating $\chi(r/L)$ gives this bound.  The coefficients
of the two uncut primitive one-forms in \eqref{eq:cartesian-seeds} are
smooth homogeneous functions of degree $-1$; exterior differentiation
therefore gives the field bounds in \eqref{eq:seed-bounds}.  The uncut
Cartesian tensor coefficients are smooth homogeneous functions of degree
$-3$, which gives the tensor bounds by the same product rule.  Smooth
cutoff jets give the estimates also at $r=L$ and $r=2L$.
\end{proof}

\subsection{Reduction to a scalar equation}

For a positive function $u$, define the physical data by
\begin{equation}
\label{eq:conformal-data}
 g=u^4\delta,\qquad K=u^{-2}\sigma,\qquad
 \mathcal E=u^{-6}\widetilde E,\qquad
 \mathcal B=u^{-6}\widetilde B.
\end{equation}
The next calculation fixes the conformal weights and the normalization
of the scalar equation directly.

\begin{lemma}[Conformal electrovacuum equations]
\label{lem:conformal}
For every smooth positive $u$, the data \eqref{eq:conformal-data} satisfy
\[
 \tau=\tr_g K=0,\qquad
 \operatorname{div}_g\mathcal E=\operatorname{div}_g\mathcal B=0,
 \qquad
 \operatorname{div}_gK=2(\mathcal E\times_g\mathcal B)^\flat_g.
\]
They satisfy the electrovacuum Hamiltonian constraint if and only if
\begin{equation}
\label{eq:scalar}
 -\Delta_\delta u=p u^{-7}+\ell u^{-3}\quad\hbox{on }\Omega,
\end{equation}
where, writing $t=\cos\theta$,
\begin{equation}
\label{eq:coefficients}
\begin{aligned}
 p&=\frac18|\sigma|_\delta^2
   =\frac{s^4c^4\sin^6\theta}{4r^6},\\
 \ell&=\frac14\bigl(|\widetilde E|_\delta^2
                  +|\widetilde B|_\delta^2\bigr)\\
   &=\frac{s^2}{4r^4}
      \left[c^2(9t^4-2t^2+1)+(rc')^2(1-t^4)\right].
\end{aligned}
\end{equation}
The nonnegative coefficients $p,\ell$ are smooth and rotation invariant.
For every Cartesian multi-index $\alpha$ of order $j\ge0$,
\begin{equation}
\label{eq:source-coefficient-bounds}
 |\partial^\alpha p|\le C_j s^4r^{-6-j}\mathbf1_{\{r\ge L\}},\qquad
 |\partial^\alpha\ell|\le C_j s^2r^{-4-j}\mathbf1_{\{r\ge L\}}.
\end{equation}
\end{lemma}

\begin{proof}
Since $\dd V_g=u^6\dd V_\delta$, the flux forms are unchanged:
\begin{equation}
\label{eq:conformal-fluxes}
 \iota_{\mathcal E}\dd V_g=\iota_{\widetilde E}\dd V_\delta,
 \qquad
 \iota_{\mathcal B}\dd V_g=\iota_{\widetilde B}\dd V_\delta.
\end{equation}
Their closedness proves the two Maxwell divergence constraints.  The
cross-product convention gives, for any vector $Z$,
\[
 (\mathcal E\times_g\mathcal B)^\flat_g(Z)
 =\dd V_g(\mathcal E,\mathcal B,Z)
 =u^{-6}\dd V_\delta(\widetilde E,\widetilde B,Z).
\]

Here is the tensor calculation in Cartesian coordinates.  Put $b=2\log u$,
so $g=e^{2b}\delta$ and $K=e^{-b}\sigma$.  The connection difference is
\[
 D^i_{jk}=\delta^i_j b_k+\delta^i_k b_j-\delta_{jk}b^i,
 \qquad b_j=\partial_jb,
\]
with indices on the right raised by $\delta$.  Symmetry and trace-freeness
give
\[
 \delta^{jk}D^l_{ki}\sigma_{lj}=0,\qquad
 \delta^{jk}D^l_{kj}=-b^l.
\]
Consequently the derivative of $e^{-b}$ cancels the second connection
contraction, yielding
\begin{equation}
\label{eq:conformal-divergence}
 \tr_gK=e^{-3b}\tr_\delta\sigma=0,\qquad
 (\operatorname{div}_gK)_i
 =e^{-3b}\delta^{jk}\partial_k\sigma_{ij}
 =u^{-6}(\operatorname{div}_\delta\sigma)_i.
\end{equation}
Together with \Cref{lem:seeds}, this is exactly the momentum constraint.

For completeness, substituting the same $D$ into
\[
 R_{ij}=\partial_kD^k_{ij}-\partial_jD^k_{ik}
          +D^k_{kl}D^l_{ij}-D^k_{jl}D^l_{ik}
\]
gives derivative terms $-b_{ij}-\delta_{ij}\Delta_\delta b$ and quadratic
terms $b_i b_j-\delta_{ij}|\nabla_\delta b|^2$.  Contracting therefore gives
\begin{equation}
\label{eq:conformal-curvature}
 R_g=e^{-2b}\bigl(-4\Delta_\delta b
                         -2|\nabla_\delta b|^2\bigr)
     =-8u^{-5}\Delta_\delta u.
\end{equation}
The norm identities are
\[
 |K|_g^2=u^{-12}|\sigma|_\delta^2,\qquad
 |\mathcal E|_g^2=u^{-8}|\widetilde E|_\delta^2,\qquad
 |\mathcal B|_g^2=u^{-8}|\widetilde B|_\delta^2.
\]
Thus the Hamiltonian constraint
$R_g-|K|_g^2=2(|\mathcal E|_g^2+|\mathcal B|_g^2)$ becomes
\eqref{eq:scalar}, with the coefficients in \eqref{eq:coefficients}.
Those explicit expressions follow from \eqref{eq:field-components} and
$|\sigma|_\delta^2=2a^2/(r^2\sin^2\theta)$.
Their smoothness, nonnegativity, and derivative bounds follow either from
their definitions as squared norms or directly from
\Cref{lem:seeds} and the product rule.
\end{proof}

We will solve \eqref{eq:scalar} with $u\to1$ at infinity and
$\partial_ru+u/2=0$ on $S$.  This boundary condition will make $S$ a future
MOTS.  A solution and the uniform estimates needed to control all other
surfaces are established next.

\section{The conformal factor}\label{sec:elliptic}

The sources begin at radius $L$.  This separation makes their effect near
$S$ of order $L^{-2}$, while their total integral is of order $L^{-1}$.
We prove both estimates by constructing the solution directly.

\begin{proposition}[Uniform conformal factor]\label{prop:conformal-factor}
There is $L_0>4$, depending only on the fixed cutoff, such that for every
$L\ge L_0$ and $s\in[0,1]$, Equation~\eqref{eq:scalar} has a smooth
axisymmetric solution $u=u_0+v$ on $\Omega$, where $u_0=1+r^{-1}$,
satisfying
\begin{equation}\label{eq:robin}
 u_r+\tfrac12u=0\quad\text{on }S,
 \qquad u\longrightarrow1\quad\text{as }r\longrightarrow\infty.
\end{equation}
It is unique among smooth solutions with $v\ge0$ and $v=O(r^{-1})$.
One may take a constant $M=3$, independent of $s,L$, such that
$1\le u_0\le u\le M$.  Uniformly in these parameters,
\begin{equation}\label{eq:radial-bounds}
 \begin{aligned}
  |v|+r|v_r|&\le Cs^2L^{-2} &&(r\ge1),\\
  |v_{rr}|&\le Cs^2L^{-2} &&(1\le r\le2).
 \end{aligned}
\end{equation}
For every integer $k\ge0$, all Cartesian derivatives satisfy
\begin{equation}\label{eq:all-derivative-bounds}
 |D^kv|\le C_k\frac{s^2}{L}r^{-1-k},
 \qquad |D^k(u-1)|\le C_k r^{-1-k}.
\end{equation}
For fixed $L$, $s\mapsto v_s$ is continuous in the sup norm, as are
\[
 s\longmapsto\int_{\R^3}H_s\,\dd x,
 \qquad
 s\longmapsto N_s(0)=\frac1{4\pi}\int_{\R^3}\frac{H_s(x)}{|x|}\,\dd x,
\]
where $H_s=pu_s^{-7}+\ell u_s^{-3}$ is extended by zero inside the unit
ball and $N_s$ is its Newton potential.  Moreover,
\begin{equation}\label{eq:source-integrals}
 0\le\int_{\R^3}H_s\,\dd x\le C\frac{s^2}{L},
 \qquad 0\le N_s(0)\le C\frac{s^2}{L^2}.
\end{equation}
At $s=0$ the solution is exactly $u_0$.
\end{proposition}

Here and below $D^k$ denotes any Cartesian derivative of order $k$;
constants may depend on this order but are independent of
$s\in[0,1]$ and $L\ge L_0$.
We first record the elementary potential estimates used in the proof.

\begin{lemma}[Newton potential and Robin reflection]\label{lem:newton}
Suppose $L>4$ and $H\in C(\R^3)$ obeys
\[
 |H(y)|\le a|y|^{-4}\mathbf1_{\{|y|\ge L\}}.
\]
Then
\[
 N_H(x)=\frac1{4\pi}\int_{\R^3}\frac{H(y)}{|x-y|}\,\dd y
\]
is well defined and $C^1$, satisfies $-\Delta N_H=H$ distributionally,
and is smooth and harmonic on $\{|x|<L\}$.  For $j=0,1$,
\begin{equation}\label{eq:newton-bounds}
 |D^jN_H(x)|\le
 \begin{cases}
  C_j aL^{-2-j},& |x|\le L/2,\\
  C_j a\bigl(r^{-2-j}+L^{-1}r^{-1-j}\bigr),&r=|x|\ge L/2.
 \end{cases}
\end{equation}
The first estimate holds for every $j\ge0$.  The operator
\[
 (\mathcal RH)(x)=N_H(x)+r^{-1}N_H(x/r^2),\qquad x\in\Omega,
\]
preserves nonnegativity, satisfies
$-\Delta\mathcal RH=H$ distributionally and
$(\partial_r+\tfrac12)\mathcal RH=0$ on $S$, and obeys
\begin{equation}\label{eq:robin-operator-bounds}
 \begin{gathered}
 \|\mathcal RH\|_\infty\le CaL^{-2},\qquad
 |\mathcal RH|+r|\partial_r\mathcal RH|\le CaL^{-2},\\
 |\partial_r^2\mathcal RH|\le CaL^{-2}\quad(1\le r\le2).
 \end{gathered}
\end{equation}
The homogeneous harmonic Robin problem on $\Omega$ has no nonzero
solution of order $O(r^{-1})$.
\end{lemma}

\begin{proof}
The bound on $H$ gives $\|H\|_{L^1}\le4\pi a/L$.  The Newton kernel
and its first derivatives are locally integrable.  Localizing around a
point and separating a small ball about its kernel singularity proves
that the potential and its first derivatives are continuous; the
singular part of the gradient integral is bounded by a constant times
the ball's radius.  The fundamental-solution identity
$-\Delta(4\pi|x|)^{-1}=\delta_0$, obtained by integration by parts around
the pole, therefore gives the distributional equation by Fubini's
theorem.

For $|x|\le L/2$, source separation allows differentiation under the
integral to every order and gives
\[
 |D^jN_H(x)|\le C_j a\int_L^\infty
          \rho^{-4}\rho^{-1-j}\rho^2\,\dd\rho
       \le C_j aL^{-2-j}.
\]
For $r=|x|\ge L/2$, split the integral at $|x-y|=r/2$.
In the inner part $|y|\ge r/2$, and for $j=0,1$ its absolute value is at
most
\[
 Ca r^{-4}\int_{|z|\le r/2}|z|^{-1-j}\,\dd z
 \le Ca r^{-2-j}.
\]
The outer part is bounded by
$Cr^{-1-j}\|H\|_{L^1}\le CaL^{-1}r^{-1-j}$.
This proves \eqref{eq:newton-bounds}.

Write $I(x)=x/r^2$.  The Kelvin identity in dimension three is
\[
 \Delta\bigl(r^{-1}N_H(I(x))\bigr)
       =r^{-5}(\Delta N_H)(I(x)).
\]
Its right side vanishes on $\Omega$ because $|I(x)|\le1<L$.
Along each ray $x=r\omega$, direct differentiation gives
\[
 (\mathcal RH)(\omega)=2N_H(\omega),\qquad
 \partial_r(\mathcal RH)(\omega)=-N_H(\omega),
\]
which proves the Robin condition including its angular dependence.
Positivity follows from the two positive kernels.  Combining
\eqref{eq:newton-bounds} in the two radial regions gives the first two
bounds in \eqref{eq:robin-operator-bounds}.  For the last bound both
potential arguments are separated from the source.  More explicitly,
put $n(t)=N_H(t\omega)$; then
\[
 \frac{\dd^2}{\dd r^2}\bigl(r^{-1}n(r^{-1})\bigr)
 =2r^{-3}n(r^{-1})+4r^{-4}n'(r^{-1})+r^{-5}n''(r^{-1}).
\]
The separated estimates for $j=0,1,2$, together with the corresponding
bound for $\partial_r^2N_H(r\omega)$, give the claim.

For uniqueness of the homogeneous problem, let $w$ be harmonic on
$\Omega$, $C^1$ up to $S$, with $w=O(r^{-1})$ and
$w_r+\tfrac12w=0$ on $S$.  Extend it to $0<r<1$ by
$r^{-1}w(I(x))$.  Values and tangential derivatives match at $S$, and
the interior radial derivative there is $-w-w_r=w_r$.
Thus the extension $\widehat w$ is distributionally harmonic across $S$
and hence harmonic on $\R^3\setminus\{0\}$. It is bounded near zero
and is $O(r^{-1})$ at infinity. For sufficiently small $\varepsilon$ and
large $R$, the maximum principle applied to $\pm\widehat w$ on
$\varepsilon<|x|<R$ gives
$|\widehat w(x)|\le C\varepsilon/|x|+C/R$, with one constant $C$:
the harmonic right side bounds the absolute value on both boundary
spheres. Letting $\varepsilon\downarrow0$ and $R\to\infty$ proves
uniqueness.
\end{proof}

\begin{lemma}[Localized derivative bootstrap]\label{lem:regularity}
Let $H$ satisfy the hypotheses of \Cref{lem:newton}.  If $H\in C^j(\R^3)$,
then $N_H\in C^{j+1}(\R^3)$.  If, in addition,
\[
 |D^iH(y)|\le a_j|y|^{-4-i}\mathbf1_{\{|y|\ge L\}}
 \quad(0\le i\le j),
\]
then \eqref{eq:newton-bounds} holds through order $j+1$, with a constant
$C_j a_j$ in place of $C_j a$.  Consequently the local gain in regularity
and the derivative decay use no first moment of $H$.
\end{lemma}

\begin{proof}
Fix a point $x_0$.  Choose a smooth compactly supported cutoff $\zeta$
equal to one near $x_0$, and keep this cutoff fixed when differentiating
with respect to $x$.  Writing $G(x)=(4\pi|x|)^{-1}$, for
$|\alpha|=j+1$ choose $\alpha=\beta+e_i$, $|\beta|=j$.  Distributional
integration by parts gives, near $x_0$,
\begin{equation}\label{eq:localized-kernel}
 D^\alpha N_H(x)
 =\int_{\R^3}\partial_iG(x-y)D_y^\beta(\zeta H)(y)\,\dd y
  +\int_{\R^3}D^\alpha G(x-y)(1-\zeta(y))H(y)\,\dd y.
\end{equation}
There are no boundary terms since $\zeta H$ has compact support.
The first integral is continuous by local integrability of
$\partial_iG$ and the same small-ball argument as above; the second is
smooth near $x_0$ by source separation and integrability at infinity.
The formula, also applied at lower orders, proves $C^{j+1}$ regularity.

For the estimate at $|x_0|=r\ge L/2$, take $\zeta$ supported in
$B(x_0,r/2)$ and equal to one on $B(x_0,r/4)$, with
$|D^i\zeta|\le C_i r^{-i}$.  On its support every Leibniz term in
$D^\beta(\zeta H)$ is bounded by $C_j a_jr^{-4-j}$.
The first integral in \eqref{eq:localized-kernel} at $x_0$ is therefore
bounded by
\[
 C_j a_jr^{-4-j}\int_{|z|\le r/2}|z|^{-2}\,\dd z
 \le C_j a_jr^{-3-j}.
\]
The second is bounded by
\[
 C_jr^{-2-j}\|H\|_{L^1}
 \le C_j a_jL^{-1}r^{-2-j}.
\]
These are precisely the claimed bounds at order $j+1$; the same
argument gives the lower orders.  On $|x|\le L/2$, differentiation of
the separated integral already gives every derivative estimate.
\end{proof}

\begin{proof}[Proof of \Cref{prop:conformal-factor}]
The coefficient estimates from \Cref{lem:seeds,lem:conformal} imply,
for every $j\ge0$,
\begin{equation}\label{eq:elliptic-coefficient-bounds}
 |D^jp|\le C_js^4r^{-6-j}\mathbf1_{\{r\ge L\}},\qquad
 |D^j\ell|\le C_js^2r^{-4-j}\mathbf1_{\{r\ge L\}}.
\end{equation}
Let $\mathcal C=\{v\in C_b(\Omega):v\ge0\}$, a complete metric space
in the sup norm.  For $v\in\mathcal C$, put
\begin{equation}\label{eq:fixed-point}
 \begin{split}
 H_v&=p(u_0+v)^{-7}+\ell(u_0+v)^{-3},\\
 N_v(x)&=\frac1{4\pi}\int_{\R^3}\frac{H_v(y)}{|x-y|}\,\dd y,\\
 T_s(v)(x)&=N_v(x)+r^{-1}N_v(x/r^2).
 \end{split}
\end{equation}
The source is extended by zero to $\R^3$; this extension is continuous
because it vanishes on $r<L$.  Since $u_0+v\ge1$, the mean-value theorem
for the negative powers gives
\begin{equation}\label{eq:source-contraction}
 \begin{split}
 0\le H_v&\le Cs^2r^{-4}\mathbf1_{\{r\ge L\}},\\
 |H_v-H_w|&\le Cs^2r^{-4}\mathbf1_{\{r\ge L\}}
                  \|v-w\|_\infty.
 \end{split}
\end{equation}
By \Cref{lem:newton}, $T_s$ maps the whole cone $\mathcal C$ into itself,
with
\[
 \|T_s(v)\|_\infty\le Cs^2L^{-2},\qquad
 \|T_s(v)-T_s(w)\|_\infty\le CL^{-2}\|v-w\|_\infty.
\]
Choose $L_0>4$ so large that the second constant is at most $1/2$ and
the first, for $s=1$, is at most one.  The contraction theorem now gives
an actual fixed point $v\in\mathcal C$, unique in that cone.
The operator commutes with rotation about the prescribed axis, so its
unique fixed point is axisymmetric.  It also gives $v=0$ when $s=0$.
\Cref{lem:newton} yields \eqref{eq:radial-bounds}, $1\le u_0\le u\le3$,
the Robin condition, and the scalar equation distributionally.
The potential bounds give $v=O(r^{-1})$, hence $u\to1$.

For completeness we bootstrap both smoothness and the full decay.
The fixed point is $C^1$ by \Cref{lem:newton}; near $S$ it is already
smooth by separation of both potential arguments from the source.
Suppose derivatives of $u-1$ through order $j$ have the uniform bounds
$C_i r^{-1-i}$.  The case $j=0$ follows from
\eqref{eq:newton-bounds} and $u_0=1+r^{-1}$.  Because $u\ge1$, the
chain and product rules with \eqref{eq:elliptic-coefficient-bounds}
give
\[
 |D^iH_v|\le C_js^2r^{-4-i}\mathbf1_{\{r\ge L\}}
 \quad(0\le i\le j).
\]
The zero extension is $C^j$ since the source vanishes near $S$.
\Cref{lem:regularity} gains one derivative for $N_v$ with the stated
bounds.  For the reflected term, the separated estimates and the
homogeneity of $I(x)$ give at every order $k$
\[
 \bigl|D^k\bigl(r^{-1}N_v(I(x))\bigr)\bigr|
       \le C_ks^2L^{-2}r^{-1-k}.
\]
In the region $1\le r\le L/2$, the ordinary term is bounded by
$C_ks^2L^{-2-k}$; in the other region it is bounded by
$C_ks^2(r^{-2-k}+L^{-1}r^{-1-k})$.
Each bound is at most $C_ks^2L^{-1}r^{-1-k}$ in its region.
This proves the next induction step and \eqref{eq:all-derivative-bounds}
at every order, with constants independent of $L,s$.
In particular $u$ is smooth and solves \eqref{eq:scalar} pointwise.
If another solution has $v\ge0$ and $v=O(r^{-1})$, subtracting
$\mathcal RH_v$ leaves a decaying homogeneous harmonic Robin solution.
\Cref{lem:newton} makes that difference zero; cone uniqueness then
identifies the solutions.

Finally, fix $L\ge L_0$.  The coefficients have the form
$p_s=s^4p_1$, $\ell_s=s^2\ell_1$.  Thus uniformly for $v\in\mathcal C$,
\[
 \|T_s(v)-T_t(v)\|_\infty\le CL^{-2}|s-t|.
\]
Using the contraction constant $1/2$ in the fixed point equations gives
\[
 \|v_s-v_t\|_\infty\le 2CL^{-2}|s-t|.
\]
Combining this with the pointwise source Lipschitz bound yields
\[
 |H_s(y)-H_t(y)|\le C|s-t|\,|y|^{-4}
                    \mathbf1_{\{|y|\ge L\}}.
\]
The right side is integrable with both weights $1$ and $|y|^{-1}$.
It proves continuity of the two source integrals, and
\eqref{eq:source-contraction} gives their quantitative bounds
\eqref{eq:source-integrals}.  Uniform convergence also gives continuity
of the restriction $u_s|_S$, which will imply continuity of the horizon
area.  No finite first moment or multipole expansion has entered the
argument.
\end{proof}

\section{The horizon and all enclosing surfaces}\label{sec:geometry}

The estimates for the conformal factor make the coordinate spheres useful
barriers.  Their positive expansion outside the boundary also gives a closed
two-form that compares the boundary area with every enclosing cut.

\begin{proposition}\label{prop:geometry}
For all sufficiently large $L$, uniformly for $0\le s\le1$, the data constructed
in \Cref{sec:construction,sec:elliptic} are smooth source-free electrovacuum
data on
\[
 \Omega=\{x\in\R^3:r\ge1\},\qquad S=\{r=1\},
\]
with $\tr_gK=0$.  The metric is complete with its boundary included, has one
asymptotically flat end, and satisfies
\begin{equation}\label{eq:geometric-falloff}
 g_{ij}-\delta_{ij}=O_4(r^{-1}),\qquad K_{ij}=O_3(r^{-3}),\qquad
 \mathcal E^i,\mathcal B^i=O_3(r^{-2}).
\end{equation}
The standard effective $U(1)$ action preserves all the data and $S$, and
$\mu$ and $|J_{\mathrm{con}}|_g$ are integrable.  With the normal toward the end,
$S$ is a future MOTS, and no compact smooth embedded enclosing surface
in $\operatorname{int}\Omega$, possibly disconnected, has everywhere
nonpositive future outward expansion.  Every enclosing cut in the sense of
\Cref{sec:introduction} has area at least $A=\Area_g(S)$.  Moreover, for a
constant $C_A$ independent of $s,L$,
\begin{equation}\label{eq:area-estimate}
 64\pi\le A\le64\pi+C_A\frac{s^2}{L^2},
 \qquad A=64\pi+O(s^2/L^2),
\end{equation}
and $A(s,L)$ is continuous in $s$ for fixed $L$.
\end{proposition}

\begin{proof}
The domain is connected and oriented by the standard orientation of $\R^3$;
its boundary is a connected two-sphere.  Outside each sufficiently large
ball it is a single product end.  Smoothness up to $S$ and across the rotation
axis follows from \Cref{lem:seeds,prop:conformal-factor}; the constraints and
maximality follow from \Cref{lem:conformal}.  The seeds are invariant under
the standard rotations, and uniqueness of the fixed point makes $u$ invariant
as well.  Thus these rotations preserve $g,K,\mathcal E,\mathcal B$, with
generator $\eta=\partial_\phi$ of period $2\pi$.  The action is effective
because a rotation fixing every point of $\Omega$ has angle $0$ modulo $2\pi$.

The bound $1\le u\le M$ gives
$\delta\le g\le M^4\delta$.  The Euclidean intrinsic length metric on the
closed exterior is complete: a Cauchy sequence has a Euclidean limit in
$\Omega$, and local smooth boundary charts show convergence also in the
intrinsic length metric.  Comparison of curve lengths now proves completeness
for $g$ with $S$ included.  The derivative bounds in
\Cref{lem:seeds,prop:conformal-factor} and the product rule give
\eqref{eq:geometric-falloff}.

The electrovacuum constraints identify
\[
 \mu=\frac{|\mathcal E|_g^2+|\mathcal B|_g^2}{8\pi}=O(r^{-4}),
 \qquad
 J_{\mathrm{con}}=\frac{(\mathcal E\times_g\mathcal B)^{\flat_g}}{4\pi}.
\]
The fields vanish for $r\le L$ and are radial, hence parallel, for $r\ge2L$.
Consequently $J_{\mathrm{con}}$ is supported in the compact annulus
$L\le r\le2L$.  Since $\dd V_g=u^6\dd x\le M^6\dd x$, both required
densities are integrable.  The horizon and enclosing-surface assertions are
proved in \Cref{lem:positive-spheres,lem:area-minimum} below.
\end{proof}

\begin{lemma}\label{lem:positive-spheres}
For $L$ sufficiently large, the coordinate sphere $S_r=\{r=\mathrm{const}\}$
has future outward expansion zero at $r=1$ and strictly positive at every
$r>1$.  No compact smooth embedded enclosing surface in
$\operatorname{int}\Omega$, with any finite number of components, is weakly
future outer trapped.
\end{lemma}

\begin{proof}
The normal to $S_r$ toward infinity is $\nu_r=u^{-2}\partial_r$.  Since
$\sigma$ has only mixed radial--azimuthal components,
$K|_{TS_r\times TS_r}=0$, including on the axis by smoothness.  The mean
curvature is therefore also the future outward expansion, and direct
divergence gives
\begin{equation}\label{eq:sphere-expansion}
 \theta_+(S_r)=H_{S_r}
 =\frac{1}{u^6r^2}\partial_r(r^2u^4)
 =\frac{2W}{ru^3},\qquad W=u+2ru_r.
\end{equation}
This formula permits angular dependence of $u$: the normal has no angular
components.  The Robin condition \eqref{eq:robin} gives $W(1,\omega)=0$.
For $r\ge2$, \eqref{eq:radial-bounds} gives
\[
 W=1-r^{-1}+v+2rv_r\ge\frac12-\frac{Cs^2}{L^2}.
\]
On $1\le r\le2$, the same estimates yield
\[
 W_r=r^{-2}+3v_r+2rv_{rr}\ge\frac14-\frac{Cs^2}{L^2}.
\]
Increasing $L$ once, uniformly over $s\in[0,1]$, we obtain
\begin{equation}\label{eq:positive-W}
 \begin{aligned}
 W_r\ge\frac18,\qquad W\ge\frac{r-1}{8}
     &\quad(1\le r\le2),\\
 W\ge\frac14&\quad(r\ge2).
 \end{aligned}
\end{equation}
Thus the strict positivity holds even arbitrarily close to $S$.

Let $\Sigma\subset\operatorname{int}\Omega$ be any compact smooth embedded
enclosing surface, without a symmetry or connectedness restriction.  To
cover also the interpretation allowing components with interior edges, let
$\Sigma_{\mathrm{cl}}$ be the union of its components without boundary.
The other components do not affect separation.  Indeed, a path avoiding
$\Sigma_{\mathrm{cl}}$ can be perturbed to miss their edges and cross their
interiors transversely at finitely many points.  For each crossing, choose
an embedded arc in that component from the crossing to an edge.  A narrow
tubular strip along the arc gives a detour on one side of the component,
around the edge, and back on the other side, removing the crossing without
introducing another.  Compactness and disjointness of the components, and
their location in $\operatorname{int}\Omega$, let every detour avoid the
other components and $S$.  Thus any path avoiding $\Sigma_{\mathrm{cl}}$
can avoid all of $\Sigma$, so $\Sigma_{\mathrm{cl}}$ still encloses $S$
and is nonempty.

Choose $p\in\Sigma_{\mathrm{cl}}$ where $r$ has its global maximum
$R>1$ on $\Sigma_{\mathrm{cl}}$.  The increasing radial ray beyond $p$
misses $\Sigma_{\mathrm{cl}}$.  Starting just beyond $p$, the same detours
past any edged components connect this side to infinity without meeting
$\Sigma$.  Hence the prescribed normal at $p$ is
$\nu_\Sigma(p)=u(p)^{-2}\partial_r$, and $T_p\Sigma=T_pS_R$.
With the convention $H_\Sigma=\operatorname{div}_\Sigma
\nu_\Sigma$, restriction of the ambient Hessian gives, at $p$,
\[
 0\ge\Delta_\Sigma r
 =\tr_{T_p\Sigma}\nabla_g^2r-H_\Sigma\,\nu_\Sigma(r)
 =|\nabla_g r|_g\bigl(H_{S_R}-H_\Sigma\bigr).
\]
For the last identity, the tangential Hessian of a defining function for a
level surface equals its second fundamental form times the length of its
gradient.  It follows that $H_\Sigma(p)\ge H_{S_R}(p)>0$.  Tangency also gives
$\tr_\Sigma K(p)=0$, so $\theta_+(\Sigma)(p)>0$.  This excludes precisely
the condition $\theta_+\le0$ everywhere.
\end{proof}

\begin{lemma}\label{lem:area-minimum}
For the same choices of $L$, every enclosing cut $\Gamma=\partial D$ from the
specified class has $\Area_g(\Gamma)\ge\Area_g(S)$.  Thus $S$ attains the
minimum enclosing area, and \eqref{eq:area-estimate} holds.
\end{lemma}

\begin{proof}
Let $\pi:\Omega\longrightarrow\Sph^2$ be the radial projection
$\pi(r,\omega)=\omega$, and let $\dd A_{\Sph^2}$ denote the unit round area
form.  Define the smooth closed two-form
\[
 \beta=\pi^*\bigl(u(1,\omega)^4\dd A_{\Sph^2}\bigr).
\]
It is closed because every two-form on $\Sph^2$ is closed.  Its comass, namely
the supremum of $|\beta(e_1,e_2)|$ over $g$-orthonormal pairs, is
\begin{equation}\label{eq:calibration-comass}
 \|\beta\|_{*,g}(r,\omega)
 =\frac{u(1,\omega)^4}{r^2u(r,\omega)^4}\le1.
\end{equation}
Indeed, $\beta$ annihilates radial vectors and its value on an oriented
orthonormal angular pair is the displayed ratio.  Furthermore,
\[
 \partial_r(ru^2)=u(u+2ru_r)=uW>0\qquad(r>1),
\]
by \Cref{lem:positive-spheres}; hence $ru(r,\omega)^2\ge u(1,\omega)^2$.
The comass estimate controls every tangent plane, so it does not require
$\Gamma$ to be a radial graph.

For completeness, take any connected smooth codimension-zero submanifold
$D\subset\Omega$ that is closed in $\Omega$, has interior contained in
$\operatorname{int}\Omega$, contains the entire sufficiently distant part
of the end, and has compact smooth embedded two-sided intrinsic boundary
$\Gamma$.  Choose $R$ beyond that boundary and beyond the radius where $D$
contains the entire end.  Then
\[
 D_R=D\cap\{r\le R\}
 \quad\text{is compact and smooth, with}\quad
 \partial D_R=\Gamma\sqcup S_R.
\]
Near $\Gamma$ the truncation leaves $D$ unchanged, including at contacts
with $S$; near $S_R$ it is a regular level-set truncation in the interior
of $D$.  These boundary pieces are disjoint, so $D_R$ has no corners.
Here $D_R$ inherits the orientation of $\Omega$, and its boundary is the
\emph{full intrinsic boundary}.  In particular, any portion coinciding with
$S$ already belongs to $\Gamma$ and contributes to this boundary integral.
Stokes' theorem with the induced boundary orientations gives
\[
 \int_\Gamma\beta=-\int_{S_R}\beta
 =-\int_{\Sph^2}u(1,\omega)^4\dd A_{\Sph^2}=-A.
\]
Using \eqref{eq:calibration-comass}, including all components, we conclude
\[
 A=\left|\int_\Gamma\beta\right|
 \le\int_\Gamma\|\beta\|_{*,g}\,\dd A_g
 \le\Area_g(\Gamma).
\]
The choice $D=\Omega$ is allowed and gives $\Gamma=S$, so the infimum equals
$A$ and is attained.  No symmetry or finite-perimeter approximation has
entered this comparison.

Finally, $u(1,\omega)=2+v(1,\omega)$ with
$0\le v(1,\omega)\le Cs^2/L^2$.  Since $u\le M$,
\[
 0\le A-64\pi
 =\int_{\Sph^2}\bigl((2+v)^4-16\bigr)\dd A_{\Sph^2}
 \le4M^3\int_{\Sph^2}v\,\dd A_{\Sph^2}
 \le C_A\frac{s^2}{L^2}.
\]
This proves \eqref{eq:area-estimate}.  Supremum-norm continuity of $v$ in $s$
from \Cref{prop:conformal-factor} also proves continuity of $A(s,L)$.
\end{proof}

\section{Asymptotic quantities and the counterexamples}
\label{sec:counterexamples}

Fix the cutoff of \Cref{sec:construction}. Throughout this section \(L\)
is above a single threshold for the preceding construction, uniformly for
\(s\in[0,1]\). Write \(H\) and \(N\) for the source and Newton potential
at the fixed point. All constants below are independent of these two
parameters once that threshold is fixed.

\subsection{Mass, momentum, and charges}

\begin{proposition}\label{prop:charges}
The constructed data have finite asymptotic quantities
\begin{equation}\label{eq:vanishing-charges}
 P_{\mathrm{ADM}}=0,\qquad Q_e=Q_b=0,
\end{equation}
and
\begin{equation}\label{eq:mass}
 m=2+2N(0)+\frac1{2\pi}\int_{\R^3}H(x)\,\dd x
   =2+O(s^2/L).
\end{equation}
In particular \(m\ge2\). If \(M\) is a uniform upper bound for \(u\), then
\begin{equation}\label{eq:mass-lower-bound}
 m-2\ge\frac{s^2}{3M^3L}.
\end{equation}
The angular momentum in \eqref{eq:adm-angular} is
\begin{equation}\label{eq:angular-momentum}
 J=-\frac{4s^2}{15}.
\end{equation}
For \(s>0\), neither electromagnetic field is identically zero.
\end{proposition}

\begin{proof}
The Cartesian decay \(K=O(r^{-3})\), with \(\tau=0\), makes the integrals
in \eqref{eq:adm-momentum} of order \(O(r^{-1})\). This proves the first
part of \eqref{eq:vanishing-charges}.

The conformal weights leave both Maxwell flux forms unchanged. On every
coordinate sphere, the electric flux is
\[
 \int_{S_r}\iota_{\mathcal E}\dd V_g
 =\int_{S_r}\dd(f\,\dd\phi)
 =2\pi\bigl(f(r,\pi)-f(r,0)\bigr)=0.
\]
The identical calculation for \(h\) gives zero magnetic flux. Equivalently,
these are integrals of globally smooth exact two-forms over closed
surfaces. The charge limits are therefore zero.

Substituting \(g_{ij}=u^4\delta_{ij}\) in \eqref{eq:adm-energy} gives
\begin{equation}\label{eq:conformal-mass-flux}
 E_{\mathrm{ADM}}
 =-\frac1{2\pi}\lim_{R\to\infty}
   \int_{S_R}u^3u_r\,\dd A_\delta.
\end{equation}
Since \(u-1=O(R^{-1})\) and \(u_r=O(R^{-2})\), replacing \(u^3\) by one
has an integrated error \(O(R^{-1})\). The source \(H\) is integrable and
the ordinary Newton potential satisfies the exact flux identity
\[
 \int_{S_R}N_r\,\dd A_\delta=-\int_{B_R}H\,\dd x.
\]
For the reflected term, smoothness of \(N\) at the origin gives
\[
 R^{-1}N(\omega/R)=\frac{N(0)}{R}+O(R^{-2}),
 \qquad
 \partial_R\bigl(R^{-1}N(\omega/R)\bigr)
       =-\frac{N(0)}{R^2}+O(R^{-3}).
\]
Together with \(u_0=1+1/r\), these identities prove the existence of the
limit and the exact formula \eqref{eq:mass}. In particular no first-moment
assumption on \(H\) or multipole expansion of the ordinary potential is
needed. The estimates of \Cref{prop:conformal-factor} give
\(N(0)=O(s^2/L^2)\), \(\int H=O(s^2/L)\), and all contributions in
\eqref{eq:mass} are nonnegative. Since the linear momentum is zero,
\(m=E_{\mathrm{ADM}}\).

On \(r\ge2L\) the cutoff is one, and
\begin{equation}\label{eq:radial-tails}
 \widetilde E=\frac{2s\cos\theta}{r^2}\partial_r,
 \qquad
 \widetilde B=\frac{s(3\cos^2\theta-1)}{r^2}\partial_r.
\end{equation}
Using just the electric field and \(u\le M\),
\[
 H\ge\ell u^{-3}\ge\frac{s^2\cos^2\theta}{M^3r^4}
 \qquad(r\ge2L).
\]
As \(\int_{\Sph^2}\cos^2\theta\,\dd A=4\pi/3\), this gives
\[
 \int H\,\dd x\ge\frac{2\pi s^2}{3M^3L}.
\]
The mass formula proves \eqref{eq:mass-lower-bound}.

Finally \(\nu_r=u^{-2}\partial_r\), \(\dd A_g=u^4\dd A_\delta\), and
\(K=u^{-2}\sigma\), so all conformal weights in the angular-momentum
integrand cancel. The gravitational flux at every coordinate sphere is
\begin{align}
 J_{\mathrm{grav}}(r)
  &:=\frac1{8\pi}\int_{S_r}K(\nu_r,\eta)\,\dd A_g\notag\\
  &=\frac1{8\pi}\int_{S_r}\sigma(\partial_r,\partial_\phi)
           \,\dd A_\delta
   =-\frac{s^2c(r)^2}{4}\int_0^\pi\sin^5\theta\,\dd\theta
   =-\frac{4s^2c(r)^2}{15}.
   \label{eq:finite-radius-angular}
\end{align}
It is constant for \(r\ge2L\), proving \eqref{eq:angular-momentum} and
existence of the prescribed limit. The two nonzero angular polynomials in
\eqref{eq:radial-tails}, multiplied by the positive factor \(u^{-6}\), also
prove the final assertion.
\end{proof}

\subsection{The conserved electromagnetic correction}

The flux \eqref{eq:finite-radius-angular} vanishes near the horizon but
not at infinity. The following calculation identifies the boundary term
responsible for the difference, using precisely our constraint signs.

\begin{proposition}\label{prop:corrected-flux}
Let \(A_B=h\,\dd\phi\), the smooth magnetic potential one-form in
\Cref{sec:construction}. Then
\begin{equation}\label{eq:corrected-divergence}
 \operatorname{div}_g\left(K(\cdot,\eta)^{\sharp_g}
                              +2h\mathcal E\right)=0.
\end{equation}
The electromagnetic boundary contribution on every coordinate sphere is
\begin{equation}\label{eq:electromagnetic-correction}
 J_{\mathrm{EM}}(r)
 :=\frac1{4\pi}\int_{S_r}h\,\iota_{\mathcal E}\dd V_g
 =\frac{4s^2c(r)^2}{15}.
\end{equation}
Consequently \(\mathcal J(r):=J_{\mathrm{grav}}(r)+J_{\mathrm{EM}}(r)=0\)
for all \(r\ge1\), whereas \(J_{\mathrm{EM}}(r)\to4s^2/15\) at infinity.
\end{proposition}

\begin{proof}
Since \(A_B\) is invariant and \(A_B(\eta)=h\), Cartan's identity gives
\[
 \iota_\eta\dd A_B=-\dd h.
\]
Using \(\dd A_B=\iota_{\mathcal B}\dd V_g\) and evaluating on
\(\mathcal E\) gives
\[
 \mathcal E(h)=-\dd V_g(\mathcal E,\mathcal B,\eta).
\]
The symmetry of \(K\) and the Killing equation imply
\[
 \operatorname{div}_g\bigl(K(\cdot,\eta)^{\sharp_g}\bigr)
   =(\operatorname{div}_gK)(\eta)
   =2\dd V_g(\mathcal E,\mathcal B,\eta).
\]
Combining this with \(\operatorname{div}_g\mathcal E=0\) proves
\eqref{eq:corrected-divergence}. Directly from the flux form,
\[
 J_{\mathrm{EM}}(r)
  =\frac12\int_0^\pi h f_\theta\,\dd\theta
  =s^2c(r)^2\int_0^\pi\sin^3\theta\cos^2\theta\,\dd\theta
  =\frac{4s^2c(r)^2}{15}.
\]
Now apply \eqref{eq:finite-radius-angular}.
\end{proof}

There is no singular gauge choice in this computation: \(A_B\) is a
globally smooth one-form. The scalar \(h=A_B(\eta)\) vanishes on the
axis, where \(\eta=0\). An additive constant \(C\) in a scalar potential
would change the boundary contribution only by \(C Q_e=0\).
The norm of \(A_B\) is \(O(r^{-1})\), while
\(A_B(\eta)=h\) need not tend to zero because \(\abs\eta=O(r)\).
Moreover, for \(s>0\), the zero-charge fields in \eqref{eq:radial-tails} retain
angularly varying \(r^{-2}\) terms. Multiplication by
\(u^{-6}=1+O(r^{-1})\) preserves these leading terms. They satisfy
\eqref{eq:falloff}, but they are not \(O(r^{-3})\).
The divergence constraint does not prohibit this: in Euclidean space,
\[
 \operatorname{div}_\delta\bigl(r^{-2}a(\omega)\partial_r\bigr)=0
\]
for every smooth angular function \(a\). No curl-free or stationarity
condition is imposed. Thus \Cref{prop:corrected-flux} is consistent with
electromagnetic angular-momentum conservation; it is the identification
of that conserved flux with \eqref{eq:adm-angular} which fails here.

\subsection{Strict violation and nondegenerate equality}

The area estimate and angular momentum already give, uniformly in the
family,
\begin{equation}\label{eq:strict-branch}
 A\ge64\pi>\frac{32\pi}{15}\ge8\pi\abs J.
\end{equation}
Since \(Q=0\), this is the strict physical area branch. Set
\begin{equation}\label{eq:defect}
 D(s,L)=m(s,L)^2-\frac{A(s,L)}{16\pi}
                         -\frac{4\pi J(s,L)^2}{A(s,L)}.
\end{equation}
By \eqref{eq:area-estimate}, \eqref{eq:mass}, and
\eqref{eq:angular-momentum},
\begin{equation}\label{eq:negative-limit}
 \lim_{L\to\infty}D(1,L)=-\frac1{225}.
\end{equation}
Hence \(D(1,L)<0\) for all sufficiently large finite \(L\).

\begin{proposition}\label{prop:equality}
For every sufficiently large fixed \(L\), there exists \(s_*\in(0,1)\)
such that \(D(s_*,L)=0\). These equality data satisfy the strict branch
\eqref{eq:strict-branch} and have no Kerr--Newman realization with the
specified asymptotic parameters and induced fields.
\end{proposition}

\begin{proof}
Write the uniform area estimate as
\[
 64\pi\le A\le64\pi+C_A s^2/L^2.
\]
The lower mass bound \eqref{eq:mass-lower-bound} and \(A\ge64\pi\) imply
\begin{equation}\label{eq:positive-small-amplitude}
 D(s,L)\ge
 \frac{4s^2}{3M^3L}-\frac{C_A s^2}{16\pi L^2}-\frac{s^4}{225}.
\end{equation}
Choose \(L\) above the thresholds for the previous propositions,
\eqref{eq:negative-limit}, and positivity of the quadratic coefficient
on the right of \eqref{eq:positive-small-amplitude}. For this fixed \(L\),
the right side is positive at sufficiently small nonzero \(s\).

The fixed point depends continuously on \(s\) in the supremum norm.
Restriction to \(S\) therefore gives continuity of \(A\), and the
integrable source bounds give continuity of \(N(0)\) and \(\int H\).
The mass formula and \eqref{eq:angular-momentum} show that \(D(s,L)\) is
continuous. Its positive value at a nonzero small parameter and its
negative value at \(s=1\) yield a root \(s_*\in(0,1)\) by the
intermediate value theorem.

At this root, \(Q_e=Q_b=0\), \(J\ne0\), and both electromagnetic fields
are nonzero. The strict branch and the equality equation give
\begin{equation}\label{eq:subextremality}
 m^2-\abs J=\frac{(A-8\pi\abs J)^2}{16\pi A}>0.
\end{equation}
Thus the Kerr--Newman member with these parameters would be a
subextremal zero-charge member, namely Kerr with zero Maxwell tensor.
Indeed the dyonic Kerr--Newman potential is linear in its electric and
magnetic charges; see \cite[Equation~(3)]{CKST2018}. When both vanish,
\(F=0\). Every spacelike slice then has
\[
 (\iota_nF)|_{T\Omega}=0,\qquad
 \star_g(F|_{T\Omega})=0,
\]
independently of the slice or its future unit normal. This contradicts
the nonzero constructed fields, and proves the asserted obstruction.
\end{proof}

\begin{proof}[Proof of \Cref{thm:main}]
\Cref{lem:seeds,lem:conformal,prop:conformal-factor,prop:geometry}
provide smooth maximal electrovacuum data with every topological,
asymptotic, boundary, outermostness, and area condition in
\Cref{sec:conventions}. \Cref{prop:charges} supplies the finite ADM and
charge limits, positive mass and zero linear momentum. The physical branch
is strict by \eqref{eq:strict-branch}. For \(s=1\) and sufficiently large
finite \(L\), \eqref{eq:negative-limit} gives the first conclusion.
For the same large \(L\), \Cref{prop:equality} gives the second.
\end{proof}

\clearpage
\bibliographystyle{alpha}
\bibliography{references/references}
\end{document}
